% =====================================================================
%  sec2.tex — 2. Mirror Symmetry And Hitchin's Equations
%  原文 pp.2–3
% =====================================================================
\section{镜像对称与 Hitchin 方程}

下一个要素是 Langlands 对偶，或称 Goddard–Nuyts–Olive 对偶。每个紧李群 $G$ 都自然地联系着一个对偶群 $\LG$：例如 $\SU(N)$ 的对偶是 $\PSU(N)=\SU(N)/\Z_{N}$，$\mathrm{E}_{8}$ 的对偶是 $\mathrm{E}_{8}$，如此等等\zhnt{2.1}。我们还须回顾 A-模型与 B-模型之间镜像对称的概念（例如见 Hori 等 (2003)）。这是二维非线性 $\sigma$ 模型的量子对称性，其最基本的作用是把复几何的问题转化为辛几何的问题。几何 Langlands 对应初看并不像镜像对称的一例，但结果它确实是。

稍有一点后见之明（这个问题最早在 Hausel 与 Thaddeus (2002) 中提出，前身是 Bershadsky 等 (1995) 与 Harvey 等 (1995) 的工作），我们可以问：$\YC(\LG,C)$ 的 B-模型是否可能镜像于 $\YC(G,C)$ 的 A-模型？即便提出了这个问题，不借助某种附加结构也难以回答它。合用的附加结构由 Hitchin 方程提供，见 Hitchin (1987a)。至此，$C$ 一直只是一个定向二维流形（紧致、无边）。但现在我们选取一个复结构，把 $C$ 看作复黎曼面。以 $G$ 为规范群的 Hitchin 方程，是关于一对 $(A,\phi)$ 的方程。这里 $A$ 是 $G$-丛 $E\to C$ 上的联络（我们要强调：$E$ 的结构群现在是紧群 $G$），$\phi$ 是 $C$ 上取值于 $\ad(E)$ 的 1-形式。Hitchin 方程模去规范群后是椭圆的，其方程组为
\begin{equation}
\begin{gathered}
F-\phi\wedge\phi=0,\\
D\phi=D\star\phi=0.
\end{gathered}
\label{eq:hitchin}
\end{equation}
这里 $\star$ 是由 $C$ 上复结构确定的 Hodge 星算子。$C$ 的复结构所起的作用就在于：它使我们能够写出最后一个方程\zhnt{2.2}。

Hitchin 方程的一个解有两种解释。一方面，给定这样一个解，我们可以定义复值联络 $\bm{A}=A+i\phi$。Hitchin 方程蕴涵相应的曲率 $\bm{F}=d\bm{A}+\bm{A}\wedge\bm{A}$ 为零，于是 Hitchin 方程的一个解定义了一个复值平坦联络，从而定义了 $\YC(G,C)$ 中的一个点\zhnt{2.3}。

另一方面，联络 $A$ 的 $(0,1)$ 部分确定了丛 $E$（更确切地说，是其复化，我们仍记作 $E$）上的一个 $\bar\partial$ 算子。在复维数为 1 的情形，$\bar\partial$ 算子上没有可积性条件，因此这个 $\bar\partial$ 算子赋 $E$ 以一个复结构，使它成为 $C$ 上的全纯 $G_{\C}$-丛。此外，记 $\phi=\varphi+\bar\varphi$，其中 $\varphi$ 与 $\bar\varphi$ 分别是 $\phi$ 的 $(1,0)$ 与 $(0,1)$ 部分。那么 Hitchin 方程蕴涵：$\varphi$ 作为 $K\otimes\ad(E)$（其中 $K$ 是 $C$ 的典范线丛）的截面是全纯的。这样的对 $(E,\varphi)$——$E\to C$ 为全纯 $G_{\C}$-丛，$\varphi\in H^{0}(C,K\otimes\ad(E))$——称为一个 Higgs 丛\zhnt{2.4}。

我们记 $\MH$ 为 Hitchin 方程之解的模空间，模去规范变换。这些方程的解可以有两种解释，这一事实意味着 $\MH$ 带有两个不同的自然复结构。在其中一个复结构（称为 $I$）下，$\MH$ 参数化半稳定 Higgs 丛 $(E,\varphi)$ 的同构类；在另一个复结构 $J$ 下，它参数化平坦 $G_{\C}$-丛的等价类，换言之，参数化同态 $\rho:\pi_{1}(C)\to G_{\C}$。$I$、$J$ 与 $K=IJ$ 合起来给出 $\MH$ 上一个自然的超凯勒结构，如 Hitchin (1987a) 所述。特别地，有全纯 2-形式 $\Omega_{I},\Omega_{J},\Omega_{K}$ 与 Kähler 形式 $\omega_{I},\omega_{J},\omega_{K}$。如在超凯勒几何中通常那样，它们由
\[
\Omega_{I}=\omega_{J}+i\omega_{K}
\]
以及这一等式的循环置换相联系\zhnt{2.5}。

在复结构 $J$ 下，$\MH$ 与我们先前描述的那个簇 $\YC$ 相同。$\YC$ 的自然全纯辛形式 $\Omega$ 就是 $i\Omega_{J}$；而定义 A-模型所用的实辛形式 $\omega=\operatorname{Im}\Omega$，与 $\omega_{K}$ 一致\zhnt{2.6}。复结构 $J$ 与全纯辛形式 $\Omega_{J}=\omega_{K}+i\omega_{I}$ 不依赖于 $C$ 上所选的复结构，这与 $\MH$ 其余的超凯勒结构形成对照\zhnt{2.7}。

\begin{remark}
顺带提一句，人们可以问：$\phi$（在当前语境中称为 Higgs 场）与粒子物理中的 Higgs 场究竟有多密切的关系？换言之，Hitchin (1987a) 引入的这套术语在多大程度上是有根据的？主要区别在于：粒子物理中的 Higgs 场是标量场，而 $\phi$ 是 $C$ 上的 1-形式（两者在各自的情形下都取值于规范群的某个表示）。不过，Hitchin 方程虽然最初是直接写下并加以研究的，却可以由 $\mathcal{N}=4$ 超对称规范理论通过一种扭转手续得到（类似于从 $\mathcal{N}=2$ 超对称规范理论得到 Donaldson 理论的那种手续）。在这一扭转手续中，$\mathcal{N}=4$ 超 Yang–Mills 理论中某些类似 Higgs 场的标量场确实被转换成了进入 Hitchin 方程的那个 Higgs 场。这为上述术语提供了一个合理的依据\zhnt{2.8}。
\end{remark}

如下文将要说明的，借助 Hitchin 方程，有可能回答 $\YC(\LG,C)$ 的 B-模型是否镜像于 $\YC(G,C)$ 的 A-模型这一问题。这一点实际上最早由 Hausel 与 Thaddeus (2002) 指出，并在 Kapustin 与 Witten (2007) 中作为理解几何 Langlands 对应的关键要素使用。
